\chapter{Rectangular Specht covolumes}

\begin{definition}[Integral orbit span in word coordinates]
\label{def:integral-specht-orbit}
\lean{BostConjecture.Specht.wordRepInt,BostConjecture.Specht.polytabloidInt,BostConjecture.Specht.orbitLattice}
\source{bost-specht-manuscript:page-0004}\leanok
Let $\mu$ be a Young diagram with $n$ boxes and $r$ rows, and let
$t:\{1,\ldots,n\}\longrightarrow\mu$ be a bijection. Set
$W=\{1,\ldots,r\}^{\{1,\ldots,n\}}$ and $P_\mathbb Z=\mathbb Z^W$.
The action is $(g x)(w)=x(w\circ g)$. Write $w_t(i)$ for the row of
$t(i)$, $\delta_w$ for a coordinate vector, and define
\[
 e_t=\sum_{h\in C_t}\operatorname{sgn}(h)h\delta_{w_t},\qquad
 L_t=\sum_{g\in\mathfrak S_n}\mathbb Z\,g e_t\subset P_\mathbb Z.
\]
This construction applies to every Young diagram and uses the same
column antisymmetrizer as the rectangular construction in the manuscript.
\end{definition}

\begin{lemma}[Integrality, finite generation, and complex span]
\label{lem:integral-specht-span}\uses{def:integral-specht-orbit}
\lean{BostConjecture.Specht.orbitLattice_fg,BostConjecture.Specht.orbitLattice_free,BostConjecture.Specht.orbitLattice_ne_bot,BostConjecture.Specht.wordRepInt_mem_orbitLattice,BostConjecture.Specht.complexCast_injective,BostConjecture.Specht.complexCast_polytabloid,BostConjecture.Specht.complex_span_orbitLattice}
\source{bost-specht-manuscript:page-0004}\leanok
The group $L_t$ is a nonzero finite free abelian group, stable under
$\mathfrak S_n$. Coefficientwise inclusion $j:P_\mathbb Z\to\mathbb C^W$
is injective, sends $e_t$ to its complex polytabloid, and satisfies
\[
 \operatorname{span}_\mathbb C j(L_t)
 =\operatorname{span}_\mathbb C\{g j(e_t):g\in\mathfrak S_n\}.
\]
\end{lemma}
\begin{proof}
All terms of $e_t$ have integer coefficients. Its finite permutation
orbit generates $L_t$. A submodule of the finite free module $\mathbb Z^W$
is free, since $\mathbb Z$ is a principal ideal domain. Multiplication
by $g$ permutes the generators, proving stability. The injective
coordinate map $\mathbb Z\hookrightarrow\mathbb C$ commutes with sums,
integer multiples and permutations, so it preserves the antisymmetrizer.
The coefficient of $w_t$ in $e_t$ is one: a column permutation preserving
$w_t$ preserves both coordinates of every box and hence is the identity.
Thus $L_t$ is nonzero. Every integral combination maps into the complex
orbit span; conversely every orbit vector belongs to $j(L_t)$. These
inclusions prove the equality of spans.
\end{proof}

\begin{lemma}[Scalar extension of the integral orbit span]
\label{lem:integral-specht-base-change}\uses{lem:integral-specht-span}
\lean{BostConjecture.Specht.complexification_injective,BostConjecture.Specht.range_complexification,BostConjecture.Specht.complexSpechtEquiv,BostConjecture.Specht.complexSpechtEquiv_tmul,BostConjecture.Specht.complexSpechtEquiv_equivariant}
\source{bost-specht-manuscript:page-0004}\leanok
The map $\mathbb C\otimes_\mathbb Z L_t\to\mathbb C^W$ given by
$z\otimes x\mapsto zj(x)$ is injective and identifies its image with
the complex Specht orbit span of $e_t$. It intertwines the position
actions on the integral scalar extension and the complex Specht module.
\end{lemma}
\begin{proof}
The torsion-free $\mathbb Z$-module $\mathbb C$ is flat. Tensoring the
inclusion $L_t\hookrightarrow\mathbb Z^W$ therefore preserves injectivity.
The coordinate isomorphism
$\mathbb C\otimes_\mathbb Z\mathbb Z^W\simeq\mathbb C^W$ sends
$z\otimes x$ to $zj(x)$. The image of the resulting injection is the
complex span of $j(L_t)$, since pure tensors span the tensor product.
The preceding lemma identifies that span. On a pure tensor, both orders
of applying a permutation and this map give $zj(gx)$, since $j$ commutes
with position permutations. Linearity extends the equality to every tensor.
\end{proof}

\begin{lemma}[Independence of the enumerating tableau]
\label{lem:integral-specht-all-tableaux}\uses{def:integral-specht-orbit}
\lean{BostConjecture.Specht.polytabloidInt_relabel,BostConjecture.Specht.orbitLattice_eq_span_polytabloids,BostConjecture.Specht.orbitLattice_eq}
\source{bost-specht-manuscript:page-0004}\leanok
For $u=t\circ g^{-1}$ one has $e_u=g e_t$. Consequently $L_t$ is
independent of $t$ and equals the integral span of all polytabloids
of shape $\mu$.
\end{lemma}
\begin{proof}
One has $w_u=w_t\circ g^{-1}$ and $C_u=g C_t g^{-1}$.
Substituting $h=g q g^{-1}$ into the antisymmetrizer gives
$\operatorname{sgn}(h)=\operatorname{sgn}(q)$ and
$h\delta_{w_u}=gq\delta_{w_t}$, hence $e_u=g e_t$.
Every enumerating tableau has this form for a unique $g$, so the set
of all polytabloids is exactly the orbit used to define $L_t$.
\end{proof}

\begin{definition}[Real Hilbert realization of the orbit lattice]
\label{def:real-specht-lattice}\uses{def:integral-specht-orbit}
\lean{BostConjecture.Specht.realCast,BostConjecture.Specht.realLattice,BostConjecture.Specht.realSpan,BostConjecture.Specht.spanLattice}
\source{bost-specht-manuscript:page-0004}\leanok
Equip $\mathbb R^W$ with the inner product
$\langle x,y\rangle=\sum_{w\in W}x_w y_w$. Let $j_\mathbb R$
be coefficientwise inclusion, $V_t=\operatorname{span}_\mathbb R
j_\mathbb R(L_t)$, and $\Lambda_t=j_\mathbb R(L_t)\subset V_t$.
The metric on $V_t$ is the restriction of this Hilbert metric.
\end{definition}

\begin{lemma}[Orthonormal word coordinates and integral discreteness]
\label{lem:real-specht-lattice}\uses{def:real-specht-lattice}
\lean{BostConjecture.Specht.realCast_injective,BostConjecture.Specht.inner_realCast,BostConjecture.Specht.realCast_single_orthonormal,BostConjecture.Specht.one_le_norm_realCast,BostConjecture.Specht.spanLattice_discreteTopology,BostConjecture.Specht.spanLattice_isZLattice,BostConjecture.Specht.latticeEquiv,BostConjecture.Specht.inner_realCast_wordRep,BostConjecture.Specht.norm_realCast_wordRep}
\source{bost-specht-manuscript:page-0004}\leanok
The vectors $j_\mathbb R(\delta_w)$ are orthonormal, and
$\langle j_\mathbb R(x),j_\mathbb R(y)\rangle=\sum_w x_w y_w$.
Every nonzero integer vector has norm at least one. The inclusion
identifies $L_t$ with $\Lambda_t$, which is a discrete full lattice
in $V_t$. Position permutations preserve the inner product and norm
of integral vectors.
\end{lemma}
\begin{proof}
The inner-product formula gives the orthonormality assertion. A nonzero
integer vector has a coordinate of absolute value at least one, and
its Hilbert norm is at least that absolute value. Differences of distinct
lattice vectors thus have norm at least one, proving discreteness.
Coefficientwise real inclusion is injective; restricting its codomain
to $V_t$ does not change its image. The real span of that image is $V_t$
by definition, so the lattice is full there. Precomposition by a position
permutation bijects the finite word set; reindexing the coordinate sum
proves preservation of inner products and hence of norms.
\end{proof}

\begin{lemma}[Integral Gram determinant and nonnegative logarithmic covolume]
\label{lem:integral-specht-covolume-nonneg}
\uses{lem:integral-specht-span,lem:real-specht-lattice}
\lean{BostConjecture.Specht.covolume_sq_eq_integer_gram,BostConjecture.Specht.one_le_covolume,BostConjecture.Specht.log_covolume_nonneg}
\source{bost-specht-manuscript:page-0004}\leanok
For any integral basis $(b_i)$ of $L_t$, put
$G_{ij}=\sum_{w\in W}(b_i)_w(b_j)_w\in\mathbb Z$. Then
\[
 \covol(\Lambda_t)^2=\det G,\qquad
 \covol(\Lambda_t)\geq1,\qquad \log\covol(\Lambda_t)\geq0.
\]
\end{lemma}
\begin{proof}
The integral basis maps to an integral basis of the full discrete
lattice $\Lambda_t$. The squared volume of its fundamental
parallelepiped equals its Gram determinant, which is exactly $\det G$
by the coordinate inner-product formula. The covolume is positive,
so this integer determinant is strictly positive and hence at least
one. Taking the nonnegative square root and then the logarithm gives
the two inequalities. A finite integral basis exists by finite freeness.
\end{proof}

\begin{lemma}[Covolume bound from an integral real basis]
\label{lem:specht-covolume-upper}\uses{lem:integral-specht-covolume-nonneg}
\lean{BostConjecture.Specht.covolume_sq_le_det_gram,BostConjecture.Specht.two_mul_log_covolume_le_log_det_gram}
\source{bost-specht-manuscript:page-0004,bost-specht-manuscript:page-0005}\leanok
Let $(v_i)_{i\in I}$ be a real basis of $V_t$ consisting of vectors
of $\Lambda_t$, and let $G=(\langle v_i,v_j\rangle)_{i,j\in I}$.
Then
\[
 1\leq\covol(\Lambda_t)^2\leq\det G,\qquad
 0\leq2\log\covol(\Lambda_t)\leq\log\det G.
\]
\end{lemma}
\begin{proof}
The subgroup $S=\sum_i\mathbb Z v_i$ is a full lattice contained
in $\Lambda_t$, because the $v_i$ are a real basis. Its covolume
is $[\Lambda_t:S]\covol(\Lambda_t)$, and the finite index is a
positive integer. Thus $\covol(\Lambda_t)\leq\covol(S)$.
The Gram formula gives $\covol(S)^2=\det G$. Both covolumes
are positive, so squaring and taking logarithms preserve the
inequalities. The lower bounds follow from the integral Gram
determinant of $\Lambda_t$.
\end{proof}

\begin{lemma}[Norm of an integral polytabloid]
\label{lem:polytabloid-norm}\uses{lem:real-specht-lattice,lem:integral-specht-all-tableaux}
\lean{BostConjecture.Specht.column_rowWord_injective,BostConjecture.Specht.polytabloid_summands_orthonormal,BostConjecture.Specht.norm_polytabloidInt_sq}
\source{bost-specht-manuscript:page-0005}\leanok
For every enumerating tableau $t$, its integral polytabloid satisfies
\[
 \|j_\mathbb R(e_t)\|^2=|C_t|.
\]
\end{lemma}
\begin{proof}
Distinct column permutations give distinct row words. Indeed, if
$w_t\circ g^{-1}=w_t\circ h^{-1}$, the permutation $h^{-1}g$
preserves the row and column of each labelled box, so it fixes every
box and is the identity. The vectors $\delta_{w_t\circ g^{-1}}$
are therefore pairwise orthogonal unit vectors. In the defining sum
for $e_t$, each occurs with coefficient $\operatorname{sgn}(g)$,
whose square is one. Pythagoras gives the claimed sum of $|C_t|$
ones.
\end{proof}

\begin{definition}[The canonical rectangular Specht lattice]
\label{def:rectangular-specht}
\uses{def:integral-specht-orbit,lem:integral-specht-all-tableaux,lem:real-specht-lattice}
\lean{BostConjecture.Specht.rectangle,BostConjecture.Specht.rectangleTableau,BostConjecture.Specht.rectangleLattice,BostConjecture.Specht.rectangleRank,BostConjecture.Specht.rectangleLogCovolume}
\source{bost-specht-manuscript:page-0004}\leanok
For $a,k\geq1$ and $d=ak$, let $M^{(k^a)}_\mathbb Z$ be the free
abelian group on row tabloids of the rectangle with $a$ rows of length
$k$. Give these tabloids the orthonormal metric. A tableau $t$ labels
its boxes by $1,\ldots,d$; its polytabloid is
$e_t=\sum_{\sigma\in C_t}\operatorname{sgn}(\sigma)\{\sigma t\}$,
where $C_t$ permutes entries within each column. Let $L_{a,d}$ be the
span of all these polytabloids, with the induced metric, and put
\[
 f_{a,d}=\rk_\mathbb Z L_{a,d},\qquad
 c_a(d)=\frac{\log\covol L_{a,d}}{f_{a,d}}.
\]
\end{definition}

\begin{lemma}[Rectangular row-word realization]
\label{lem:rectangular-row-words}\uses{def:rectangular-specht,lem:real-specht-lattice}
\lean{BostConjecture.Specht.colLen_rectangle,BostConjecture.Specht.rowLen_rectangle,BostConjecture.Specht.rowWord_hasRectangleContent,BostConjecture.Specht.exists_tableau_rowWord_iff,BostConjecture.Specht.mem_rectangle_rowWord_orbit_iff,BostConjecture.Specht.rectangle_orbitLattice_supported}
\source{bost-specht-manuscript:page-0004}\leanok
For $k>0$, the rectangle has $a$ rows. Its row words, and the position
permutation orbit of any one of them, are exactly the words in $a$
letters in which each letter occurs $k$ times. Every element of
$L_{a,d}$ has zero coordinates outside this set. Thus its word metric
is the orthonormal row-tabloid metric.
\end{lemma}
\begin{proof}
The rectangle is the set of pairs $(i,j)$ with $0\leq i<a$ and
$0\leq j<k$. In a tableau, entries in row $i$ correspond bijectively
to its $k$ columns. Conversely, for a word with the stated content,
choose a bijection between each letter's fiber and the $k$ columns.
These bijections together give an enumerating tableau with that row
word. Any two enumerating tableaux differ by a permutation of the
entries, so all such words form one position-permutation orbit.
Every summand of every polytabloid is one of these row words, proving
the support assertion. A row tabloid is exactly the assignment of
each entry to a row, and distinct word coordinate vectors are
orthonormal; this proves the metric assertion.
\end{proof}

\begin{lemma}[Rectangular column count and initial norm]
\label{lem:rectangular-initial-norm}\uses{def:rectangular-specht,lem:polytabloid-norm}
\lean{BostConjecture.Specht.rectangleColumnGroupEquiv,BostConjecture.Specht.card_rectangle_columnGroup,BostConjecture.Specht.norm_rectangle_polytabloidInt_sq}
\source{bost-specht-manuscript:page-0005}\leanok
For every tableau of the $a$ by $k$ rectangle,
$|C_t|=(a!)^k$ and $\|e_t\|^2=(a!)^k$.
\end{lemma}
\begin{proof}
Restricting a column permutation to the row indices in each column
gives a bijection $C_t\simeq\prod_{j=1}^k\mathfrak S_a$.
Conversely independent row permutations within each column act
on the entries through the enumerating tableau. These operations
are inverse, so the cardinality is $(a!)^k$. The preceding
polytabloid norm lemma gives the asserted squared norm.
\end{proof}

\begin{lemma}[Rectangular polytabloids as alternating columns]
\label{lem:rectangular-column-tensor}
\lean{BostConjecture.Specht.rectangleColumnGroupOf_sign,BostConjecture.Specht.polytabloidInt_rectangle_eq_prod}
\uses{lem:rectangular-row-words,lem:rectangular-initial-norm}
\source{bost-specht-manuscript:page-0008,bost-specht-manuscript:page-0009}\leanok
Let $k>0$ and let $t$ enumerate the $a$ by $k$ rectangle. Write
$p_t(i,j)$ for the entry in row $i$, column $j$. For a word $w$ in
$a$ letters, let $w_j(i)=w(p_t(i,j))$. Define
$\epsilon(f)=\operatorname{sgn}(f)$ if $f$ is a permutation of the
$a$ letters, and $\epsilon(f)=0$ otherwise. Then
\[
 e_t(w)=\prod_{j=1}^k\epsilon(w_j).
\]
\end{lemma}
\begin{proof}
Every column permutation is uniquely a tuple $(\sigma_j)_{j=1}^k$
of permutations of the rows. Its sign is $\prod_j\operatorname{sgn}(\sigma_j)$:
the column permutations have disjoint supports, so their signs multiply,
and transporting a permutation along a bijection preserves its sign.
Its contribution to $e_t(w)$ is nonzero exactly when
$w_j=\sigma_j^{-1}$ for every $j$. Such a tuple exists if and only
if all $w_j$ are permutations, and then it is unique. Its sign is
$\prod_j\operatorname{sgn}(w_j)$, since inversion preserves sign.
If any $w_j$ is not a permutation, both sides vanish.
\end{proof}

\begin{lemma}[Determinant covariance in word coordinates]
\label{lem:rectangular-word-covariance}
\lean{BostConjecture.Specht.wordTensorAction,BostConjecture.Specht.wordTensorAction_polytabloid,BostConjecture.Specht.wordTensorAction_of_mem_orbitLattice,BostConjecture.Specht.wordTensorAction_of_mem_span}
\uses{lem:rectangular-column-tensor,lem:integral-specht-all-tableaux}
\source{bost-specht-manuscript:page-0008,bost-specht-manuscript:page-0009}\leanok
For a commutative ring $R$ and an $a$ by $a$ matrix $u$, define the
matrix action on word coefficients by
\[
 (T_u x)(w)=\sum_v\left(\prod_{q=1}^{ak}u_{w(q),v(q)}\right)x(v).
\]
For every rectangular polytabloid, and hence every vector in the
$R$-span of the coefficientwise image of $L_{a,ak}$, one has
$T_u x=(\det u)^k x$. No invertibility or characteristic assumption
is required.
\end{lemma}
\begin{proof}
For one column and a function $f$ on its row indices, expansion of
the determinant gives
\[
 \sum_g\prod_i u_{f(i),g(i)}\epsilon(g)
 =\det(u_{f(i),j})_{i,j}=\epsilon(f)\det u.
\]
If $f$ is bijective this is the row-permutation rule for determinants;
otherwise two rows coincide and both sides are zero. The word sum
for a rectangular polytabloid factors into $k$ independent column
sums by distributivity. Applying the one-column identity to each
factor and the preceding coefficient formula proves
$T_u e_t=(\det u)^k e_t$. Every lattice vector is an integer linear
combination of polytabloids, so linearity proves the assertion on
the stated scalar span.
\end{proof}

\begin{lemma}[Positive rectangular rank and nonnegative covolume slope]
\label{lem:rectangular-covolume-nonneg}\uses{def:rectangular-specht,lem:integral-specht-covolume-nonneg}
\lean{BostConjecture.Specht.rectangleRank_pos,BostConjecture.Specht.rectangleLogCovolume_nonneg}
\source{bost-specht-manuscript:page-0004}\leanok
The rank $f_{a,d}$ is positive and $c_a(d)\geq0$.
\end{lemma}
\begin{proof}
The row-word coefficient of a generating polytabloid is one,
so the finite free lattice is nonzero and has positive rank.
Its covolume is at least one. Its logarithm is therefore nonnegative,
and division by the positive rank preserves that inequality.
\end{proof}

\begin{definition}[Standard enumerating tableaux]
\label{def:standard-tableau}\uses{def:integral-specht-orbit}
\lean{BostConjecture.Specht.RowStandard,BostConjecture.Specht.ColumnStandard,BostConjecture.Specht.IsStandard}
\source{bost-specht-manuscript:page-0004}\leanok
An enumerating tableau is row-standard if the labels increase from
left to right in each row, and column-standard if they increase from
top to bottom in each column. It is standard if both conditions hold.
Order row words lexicographically in increasing order of the labels.
\end{definition}

\begin{lemma}[Row-standard tableaux are determined by their row words]
\label{lem:row-standard-injective}\uses{def:standard-tableau}
\lean{BostConjecture.Specht.RowStandard.eq_of_rowWord_eq}
\source{bost-specht-manuscript:page-0004}\leanok
Two row-standard tableaux of the same shape with the same row word
are equal.
\end{lemma}
\begin{proof}
Induct on the label $i$, assuming all earlier labelled boxes agree.
The boxes of $i$ in the two tableaux have the same row. If the first
lies to the right of the second, the label $j$ occupying the latter
box in the first tableau satisfies $j<i$ by row-standardness.
By the induction hypothesis that box also contains $j$ in the second
tableau, where it contains $i$, a contradiction. Interchanging the
tableaux excludes the opposite strict column inequality. Thus the
boxes of $i$ agree, completing the induction.
\end{proof}

\begin{lemma}[Leading row word of a column-standard polytabloid]
\label{lem:column-standard-leading-word}\uses{def:standard-tableau,lem:integral-specht-span}
\lean{BostConjecture.Specht.ColumnStandard.rowWord_lt,BostConjecture.Specht.ColumnStandard.rowWord_le_of_polytabloid_ne_zero}
\source{bost-specht-manuscript:page-0004}\leanok
If $t$ is column-standard and $1\ne g\in C_t$, then
$w_t<w_t\circ g$ lexicographically. Consequently every nonzero
coordinate of $e_t$ occurs at a word at least $w_t$.
\end{lemma}
\begin{proof}
Choose the least label $i$ moved by $g$. Every smaller label is fixed.
One has $g(i)>i$: otherwise $g(i)<i$ is fixed, contradicting
injectivity of $g$. Since $i$ and $g(i)$ occupy the same column,
column-standardness forces the latter to lie in a strictly lower
row. Hence the first difference between $w_t$ and $w_t\circ g$
is an increase at $i$. Every coordinate in the polytabloid sum is
of the form $w_t\circ g^{-1}$ with $g\in C_t$, proving the second
claim, including the identity permutation.
\end{proof}

\begin{theorem}[Independence of standard polytabloids]
\label{lem:standard-polytabloid-independent}
\uses{lem:row-standard-injective,lem:column-standard-leading-word,lem:integral-specht-span}
\lean{BostConjecture.Specht.linearIndependent_standard_polytabloid}
\source{bost-specht-manuscript:page-0004}\leanok
For any Young diagram and any ring $R$, the standard integral
polytabloids, mapped coefficientwise into $R$, are linearly independent
over $R$.
\end{theorem}
\begin{proof}
Suppose a finite linear relation has a nonzero coefficient. Among
the tableaux with nonzero coefficient, choose $t$ with least row word.
Evaluate the relation at $w_t$. If another tableau $u$ contributed,
its leading-word bound would give $w_u\leq w_t$, while minimality
gives $w_t\leq w_u$. Equality of the row words implies $u=t$ by
row-standard injectivity. The coefficient of $w_t$ in $e_t$ is one,
so the evaluated relation is precisely the nonzero coefficient of
$t$, a contradiction. No division is used, so the argument holds
over every ring.
\end{proof}

\begin{lemma}[Garnir cancellation]
\label{lem:garnir-cancellation}\uses{lem:integral-specht-span}
\source{bost-specht-manuscript:page-0004}\lean{BostConjecture.Specht.sum_sign_smul_polytabloidInt_relabel_eq_zero,BostConjecture.Specht.polytabloidRat_mem_span_relabel}\leanok
Let $t$ enumerate a Young diagram and let $X$ be a set of labels.
Suppose every label of $X$ lies in a column of length at most $r$,
and $|X|>r$. Write $G_X$ for the permutations supported on $X$.
Then
\[
 \sum_{\sigma\in G_X}\operatorname{sgn}(\sigma)e_{\sigma t}=0.
\]
Moreover, $e_t$ belongs to the rational span of those
$e_{\sigma t}$ with $\sigma\in G_X\setminus C_t$.
\end{lemma}
\begin{proof}
Expand $e_t$ as its signed sum of row words obtained by column
permutations. In each such word the labels of $X$ occupy at most
$r$ rows, so two distinct labels of $X$ have the same row.
Their transposition $\tau$ fixes this word. Pairing $\sigma$ with
$\sigma\tau$ in its antisymmetrization over $G_X$ cancels the
terms, since the signs are opposite. Linearity gives the identity.
If $\sigma\in C_t$, reindexing the column sum gives
$e_{\sigma t}=\operatorname{sgn}(\sigma)e_t$. Thus the internal
terms contribute $Ne_t$, where $N=|G_X\cap C_t|\geq1$.
Moving the remaining terms to the other side and dividing by
$N$ proves the rational span assertion.
\end{proof}

\begin{lemma}[Garnir set geometry]
\label{lem:garnir-set}\uses{def:standard-tableau}
\source{bost-specht-manuscript:page-0004}\lean{BostConjecture.Specht.card_garnirSet,BostConjecture.Specht.colLen_colIndex_le_colLen_of_mem_garnirSet,BostConjecture.Specht.mem_columnGroup_iff_image_garnirSetLeft_eq}\leanok
For a cell $(i,j+1)$, let $A$ consist of the labels in column $j$
weakly below row $i$, and $B$ of those in column $j+1$ weakly above
row $i$. Then $X=A\sqcup B$ has $\mu'_j+1$ labels, all in columns
of length at most $\mu'_j$. A permutation supported on $X$
preserves columns exactly when it preserves $A$.
\end{lemma}
\begin{proof}
The two cardinalities are $\mu'_j-i$ and $i+1$; both columns
reach row $i$. Columns shorten to the right. On $X$, the column
is $j$ on $A$ and $j+1$ on $B$, so preserving columns is equivalent
to preserving $A$.
\end{proof}

\begin{theorem}[Rational standard-polytabloid spanning]
\label{lem:rational-standard-spanning}
\uses{def:standard-tableau,lem:garnir-cancellation,lem:garnir-set}
\source{bost-specht-manuscript:page-0004}\lean{BostConjecture.Specht.polytabloidRat_mem_span_standard}\leanok
Every polytabloid of a Young diagram is a rational linear combination
of standard polytabloids of the same shape.
\end{theorem}
\begin{proof}
Write $C(t)=\sum_{\ell=0}^{n-1}\operatorname{col}_t(\ell)\ell$,
$R(t)=\sum_{\ell=0}^{n-1}\operatorname{row}_t(\ell)\ell$
and $W(t)=(n^3+1)C(t)+R(t)$. Both moments are at most $n^3$,
since there are $n$ labels and each factor is at most $n$.
Consequently $W$ is bounded above by $(n^3+1)n^3+n^3$.
Induct on this bound minus $W(t)$.
If a column is not increasing, swap a larger label $x$ in row
$i$ with a smaller label $y$ in row $i'>i$ of that column.
This preserves $C$ and increases $R$ by $(i'-i)(x-y)>0$.
The polytabloid changes only by sign, so induction applies.
Otherwise the columns increase. If a row fails to increase, some
adjacent pair in that row is inverted: otherwise transitivity of
all adjacent inequalities orders the row. Apply Garnir cancellation
at this pair. It expresses $e_t$ as a rational combination of
polytabloids indexed by permutations of $X=A\sqcup B$ which fail
to preserve $A$. Column-standardness and the chosen inversion
imply that every label of $A$ exceeds every label of $B$.
Under relabeling, $C(\sigma t)=\sum_\ell
\operatorname{col}_t(\ell)\sigma(\ell)$. The contribution outside
$X$ is fixed, and that on $X$ is
$j\sum_{\ell\in X}\sigma(\ell)+\sum_{\ell\in B}\sigma(\ell)$.
The first sum is fixed. If $A$ is not preserved, neither is $B$.
The equally sized sets $\sigma(B)\setminus B$ and
$B\setminus\sigma(B)$ are nonempty; the first lies in $A$ and
its labels exceed every label in the second. Their sums are
therefore strictly ordered, which proves the assertion.
Their column moments therefore strictly increase. Each such
increase raises $W$, even if the row moment drops by as much as
$n^3$. Induction again applies. If neither case occurs, $t$ is
standard and its polytabloid is already a generator.
\end{proof}

\begin{remark}
This rational straightening proof follows the Garnir and bounded-moment
argument in the pinned TauCeti formalization. It uses the existing word
realization and divides only by the positive integer counting internal
permutations. It makes no assertion of integral standard spanning.
\end{remark}

\begin{lemma}[Real standard basis and integral rank]
\label{lem:real-standard-basis}
\uses{lem:rational-standard-spanning,lem:standard-polytabloid-independent,lem:integral-specht-all-tableaux,lem:real-specht-lattice}
\lean{BostConjecture.Specht.realSpan_eq_span_standard,BostConjecture.Specht.standardRealBasis,BostConjecture.Specht.standardRealBasis_coe,BostConjecture.Specht.standardRealBasis_mem_spanLattice,BostConjecture.Specht.finrank_orbitLattice_eq_card_standard}
\source{bost-specht-manuscript:page-0004}\leanok
For every Young diagram, the standard polytabloids form a real basis
of $V_t=\operatorname{span}_{\mathbb R}j_{\mathbb R}(L_t)$,
and every basis vector belongs to $\Lambda_t$.
Their number equals $\operatorname{rk}_{\mathbb Z}L_t$;
in particular for a rectangle it equals $f_{a,d}$.
\end{lemma}
\begin{proof}
Rational spanning expresses each polytabloid as a rational linear
combination of standard polytabloids. Coefficientwise inclusion
$\mathbb Q\hookrightarrow\mathbb R$ preserves that combination,
so each orbit generator lies in the real span of the standard family.
Integral combinations of orbit generators do also. Conversely each
standard polytabloid belongs to $L_t$, giving equality with $V_t$.
The independence theorem with $R=\mathbb R$ proves independence in
real coordinates and hence in the Hilbert realization. These vectors
therefore form a real basis, and they lie in $\Lambda_t$ by construction.
The rank of the full discrete lattice $\Lambda_t$ equals the dimension
of $V_t$, hence the number of standard tableaux. The integral
isomorphism $L_t\simeq\Lambda_t$ gives the stated rank formula.
\end{proof}

\begin{lemma}[Standard Gram bound for the orbit lattice]
\label{lem:standard-gram-covolume}
\uses{lem:real-standard-basis,lem:specht-covolume-upper}
\lean{BostConjecture.Specht.covolume_sq_le_det_standard_gram,BostConjecture.Specht.two_mul_log_covolume_le_log_det_standard_gram}
\source{bost-specht-manuscript:page-0004,bost-specht-manuscript:page-0005}\leanok
For the standard polytabloid family of any Young diagram,
\[
 \covol(\Lambda_t)^2\leq\det\Gram((e_u)_{u\in\SYT(\mu)}),\qquad
 2\log\covol(\Lambda_t)\leq
 \log\det\Gram((e_u)_{u\in\SYT(\mu)}).
\]
\end{lemma}
\begin{proof}
The preceding lemma supplies a real basis whose vectors belong to
$\Lambda_t$. Apply the integral-real-basis covolume bound to it.
The coordinate identity for this basis identifies its Gram matrix
with that of the standard polytabloids in the word metric.
\end{proof}

\begin{theorem}[Absolute irreducibility of the rectangular Specht lattice]
\label{lem:rectangular-absolute-irreducible}
\uses{def:rectangular-specht,lem:integral-specht-base-change}
\source{bost-specht-manuscript:page-0004}\notready
The characteristic-zero representation of $\mathfrak S_d$ on
$L_{a,d}$ is absolutely irreducible.
\end{theorem}

\begin{theorem}[Rectangular invariant-space identification]
\label{lem:rectangular-sl-invariants}\uses{def:rectangular-specht}
\source{bost-specht-manuscript:page-0004}\notready
For every $a$-dimensional vector space $A$ over a characteristic-zero
field, the image of $(\det A)^{\otimes k}\otimes L_{a,d}$ in
$A^{\otimes d}$ is exactly $(A^{\otimes d})^{\SL(A)}$.
\end{theorem}

\begin{theorem}[Rectangular Specht realization]
\label{lem:rectangular-invariants}
\uses{lem:rectangular-absolute-irreducible,lem:rectangular-sl-invariants}
\source{bost-specht-manuscript:page-0004}\notready
The characteristic-zero representation of $\mathfrak S_d$ on $L_{a,d}$ is
absolutely irreducible. For every $a$-dimensional vector space $A$
over a characteristic-zero field, the image of
$(\det A)^{\otimes k}\otimes L_{a,d}$ in $A^{\otimes d}$ is exactly
$(A^{\otimes d})^{\SL(A)}$.
\end{theorem}
\begin{proof}
Apply the absolute irreducibility theorem and the invariant-space
identification to the same lattice.
\end{proof}

\begin{remark}
The manuscript defines its lattice by standard polytabloids. Here
$L_{a,d}$ is the all-polytabloid orbit lattice, and $S_{a,d}$ denotes
the manuscript's standard-polytabloid lattice. The argument uses
$S_{a,d}\subseteq L_{a,d}$ and their equal rational spans; it does
not require equality over the integers. The Gram identity below
belongs to $S_{a,d}$ and gives an upper bound for the covolume of
$L_{a,d}$. Integral contraction uses $L_{a,d}$ itself, and its degree
is computed using any integral basis of that lattice. Its image
need not be the saturated invariant lattice. Determinant covariance,
absolute irreducibility and the full invariant generic fiber remain
separate properties of the all-polytabloid construction.
\end{remark}

\begin{lemma}[Monic orthogonalization of standard polytabloids]
\label{lem:orthogonal-standard-basis}\uses{lem:real-standard-basis,lem:standard-gram-covolume}
\lean{BostConjecture.Specht.orthogonalStandardBasis,BostConjecture.Specht.orthogonalStandardBasis_eq_sub_sum,BostConjecture.Specht.orthogonalStandardBasis_orthogonal,BostConjecture.Specht.orthogonalStandardBasis_triangular,BostConjecture.Specht.orthogonalStandardBasis_monic,BostConjecture.Specht.det_orthogonalStandardBasis,BostConjecture.Specht.det_standard_gram_eq_prod_orthogonal_norm_sq,BostConjecture.Specht.two_mul_log_covolume_le_sum_log_orthogonal_norm_sq}
\source{ryom-hansen-seminormal-denominators-v1}\leanok
For any total ordering of the standard tableaux, define recursively
\[
 f_t=e_t-\sum_{s<t}\frac{\langle f_s,e_t\rangle}{\|f_s\|^2}f_s.
\]
This is an orthogonal basis of $V_t$. Its change-of-basis matrix from
the standard polytabloids is triangular with diagonal entries one, and
\[
 \det\Gram((e_t))=\prod_t\|f_t\|^2,\qquad
 2\log\covol\Lambda_t\leq\sum_t\log\|f_t\|^2.
\]
\end{lemma}
\begin{proof}
Induction on the ordering shows that $f_t-e_t$ belongs to the span of
the earlier standard polytabloids, and that the earlier $f_s$ span the
same subspace. Thus $f_t$ is nonzero by independence of the $e_t$.
For $s<t$, taking the inner product with $f_s$ in the displayed
recursion cancels exactly the summand indexed by $s$; all other
summands vanish by inductive orthogonality. The resulting triangular
matrix has diagonal entries one, hence determinant one. Transformation
of Gram matrices preserves their determinant under this change of
basis. The new Gram matrix is diagonal with positive entries
$\|f_t\|^2$, giving the product formula. Take logarithms and apply
the standard Gram covolume bound.
\end{proof}

\begin{lemma}[Real isometric permutation action]
\label{lem:real-specht-action}\uses{lem:integral-specht-all-tableaux,lem:real-specht-lattice}
\lean{BostConjecture.Specht.wordRepReal,BostConjecture.Specht.realSpechtRep,BostConjecture.Specht.inner_realSpechtRep,BostConjecture.Specht.realSpechtRep_swap_symmetric,BostConjecture.Specht.realSpechtRep_standardRealBasis_coe}
\source{bost-specht-manuscript:page-0004}\leanok
The position action $P_gx(w)=x(w\circ g)$ preserves the real
Specht span and its word inner product. It agrees with relabeling
on polytabloids. Every transposition acts by a symmetric operator.
\end{lemma}
\begin{proof}
Composition of position permutations gives $P_gP_h=P_{gh}$.
The map $w\mapsto w\circ g$ is a bijection on words; reindexing
the finite sum defining the inner product proves
$\langle P_gx,P_gy\rangle=\langle x,y\rangle$.
The integral lattice is stable under $P_g$, and taking its real
span preserves stability. The integral relabeling identity extends
coefficientwise to real vectors. Applying the inner-product identity
to $x$ and $P_{g^{-1}}y$ gives
$\langle P_gx,y\rangle=\langle x,P_{g^{-1}}y\rangle$.
For a transposition $g^{-1}=g$, proving symmetry.
\end{proof}

\begin{lemma}[Symmetric commuting Murphy operators]
\label{lem:murphy-operators}\uses{lem:real-specht-action}
\lean{BostConjecture.Specht.murphyOperator,BostConjecture.Specht.murphyOperator_symmetric,BostConjecture.Specht.murphyOperator_commute}\leanok
On the real Specht span, with labels $0,\ldots,n-1$, put
$J_j=\sum_{i<j}P_{(i,j)}$. Each $J_j$ is symmetric and the
operators $J_j$ commute pairwise.
\end{lemma}
\begin{proof}
Symmetry follows by summing the symmetry of the transpositions.
If a permutation $g$ fixes $j$ and permutes the labels below $j$,
then conjugation sends $(i,j)$ to $(g(i),j)$ and hence fixes $J_j$
after reindexing its sum. For $h<j$ every transposition $(i,h)$
with $i<h$ has this property. Consequently every summand of $J_h$
commutes with $J_j$, and so does $J_h$. Equal indices commute
trivially; reversed indices follow by symmetry of commutation.
\end{proof}

\begin{lemma}[Adjacent Murphy recurrence]
\label{lem:murphy-recurrence}\uses{lem:murphy-operators}
\lean{BostConjecture.Specht.murphyOperator_succ,BostConjecture.Specht.murphyOperator_mul_swap}\leanok
For $s=P_{(i,i+1)}$, one has $J_{i+1}=sJ_i s+s$ and
$J_{i+1}s=sJ_i+1$.
\end{lemma}
\begin{proof}
Conjugation by $(i,i+1)$ sends each $(k,i)$ with $k<i$ to
$(k,i+1)$. Their sum is $J_{i+1}$ minus its last summand $s$,
proving the first identity. Multiply on the right by $s$ and use
$s^2=1$ to obtain the second.
\end{proof}

\begin{lemma}[Content coefficient at the leading row word]
\label{lem:murphy-content-coefficient}
\lean{BostConjecture.Specht.polytabloidInt_rowWord_swap,BostConjecture.Specht.sum_polytabloidInt_rowWord_swap,BostConjecture.Specht.murphyOperator_standardRealBasis_rowWord}\leanok
\uses{def:standard-tableau,lem:integral-specht-span,lem:murphy-operators}
For distinct labels $i,j$ in any tableau $t$,
\[
 e_t(w_t\circ(i,j))
 =\mathbf1_{\operatorname{row}_t(i)=\operatorname{row}_t(j)}
  -\mathbf1_{\operatorname{col}_t(i)=\operatorname{col}_t(j)}.
\]
Consequently, for a standard tableau $t$ and zero-based label $j$,
\[
 (J_je_t)(w_t)=\operatorname{col}_t(j)-\operatorname{row}_t(j).
\]
This is a word coefficient, without an assertion about the other
coefficients in the standard-polytabloid expansion.
\end{lemma}
\begin{proof}
If the two labels share a row, the swapped row word is unchanged and
its coefficient is one. If they share a column, the transposition
belongs to the column group; column antisymmetry gives coefficient
minus one. Distinct labels cannot share both coordinates.
In the remaining case suppose a column permutation $g$ produces the
swapped row word. Every label other than $i,j$ has unchanged row and
column and is therefore fixed by $g$. Since $g$ is bijective, it either
fixes both remaining labels or swaps them. The former forces their
rows equal and the latter their columns equal, both contradictions.
There is thus no term contributing to the coefficient, which is zero.
Summing over $i<j$ counts the earlier labels in the row of $j$ minus
those in its column. Standardness identifies those sets respectively
with the boxes strictly to the left and strictly above $j$. Their
cardinalities are its zero-based column and row coordinates.
\end{proof}

\begin{lemma}[Column antisymmetrizer on the tabloid module]
\label{lem:column-antisymmetrizer-line}
\uses{def:integral-specht-orbit}
\lean{OAI.Saxl.columnAlt_tabloid}\leanok
Let $P_t$ be the complex span of the permutation orbit of
$\delta_{w_t}$ and let $A_t=\sum_{g\in C_t}\operatorname{sgn}(g)P_g$.
Then $A_t(P_t)\subseteq\mathbb C e_t$.
\end{lemma}
\begin{proof}
It suffices to consider a word in the orbit of $w_t$. If two labels
in a column have the same letter, their transposition fixes the word;
pairing terms of the antisymmetrizer cancels them. Otherwise each
column has distinct letters. The number of occurrences of letter $r$
is the length of row $r$. Letter $0$ therefore occurs in every nonempty
column: there are exactly as many such columns as its multiplicity,
and each contains it at most once. Remove its occurrence from each
column. Inductively, after removing letters below $r$, the remaining
columns are exactly those of height greater than $r$, whose number is
the multiplicity of $r$. Thus each of these columns contains letter
$r$. Consequently the letters in each column are exactly its original
row indices. A unique column permutation carries the original row
word to this word. Reindexing the antisymmetrizer then gives its sign
times $e_t$. In either case the image belongs to the stated line;
linearity completes the proof.
\end{proof}

\begin{lemma}[Central transposition sum acts by total content]
\label{lem:central-content-action}
\lean{BostConjecture.Specht.diagramContent,BostConjecture.Specht.centralWordOperator_polytabloid,BostConjecture.Specht.sum_murphyOperator_eq,BostConjecture.Specht.sum_wordRepInt_eq_content_of_mem_orbitLattice}\leanok
\uses{lem:murphy-content-coefficient,lem:column-antisymmetrizer-line,lem:real-standard-basis,lem:real-specht-action}
Let $C_\mu=\sum_{(r,c)\in\mu}(c-r)$ and $Z=\sum_jJ_j$.
On the integral polytabloid lattice and its real span, $Zx=C_\mu x$.
\end{lemma}
\begin{proof}
The sum $Z$ contains each transposition once. Conjugation by any
permutation permutes these transpositions, so $Z$ commutes with the
position action. In complex word coordinates write
$e_t=A_t\delta_{w_t}$, where $A_t$ is the column antisymmetrizer.
Then $Ze_t=A_tZ\delta_{w_t}$. The vector $Z\delta_{w_t}$ lies in
the row-tabloid module. The column antisymmetrizer has image contained
in $\mathbb C e_t$ on that module, by the Ferrers column argument
for Specht irreducibility. Thus $Ze_t=\alpha e_t$ for some scalar.
For a standard $t$, evaluate at $w_t$: the preceding coefficient
formula gives $\alpha=\sum_j\operatorname{content}_t(j)=C_\mu$.
Standard polytabloids are a real basis, so linearity proves the
identity on the real span. Every integral vector embeds in that
span, and coefficientwise inclusion is injective, giving the
integral identity as well.
\end{proof}

\begin{lemma}[Standard prefix diagrams and their contents]
\label{lem:standard-prefix-diagram}\uses{def:standard-tableau,lem:central-content-action}
\lean{BostConjecture.Specht.prefixShape,BostConjecture.Specht.prefixTableau,BostConjecture.Specht.prefixTableau_isStandard,BostConjecture.Specht.card_prefixShape,BostConjecture.Specht.prefixShape_mono,BostConjecture.Specht.diagramContent_prefixShape,BostConjecture.Specht.diagramContent_prefixShape_succ}\leanok
For a standard tableau $t$ on $n$ boxes and $0\leq m\leq n$, the
boxes carrying labels $0,\ldots,m-1$ form a Young diagram
$\lambda_t^{(m)}$ of cardinality $m$. Restricting $t$ to these labels
is a standard enumerating tableau of this diagram. These diagrams
increase with $m$, and
\[
 C_{\lambda_t^{(m)}}=\sum_{j<m}\operatorname{content}_t(j).
\]
For $m<n$, the difference between the contents at $m+1$ and $m$
is the content of the box carrying $m$.
\end{lemma}
\begin{proof}
The label in a box weakly northwest of a second box is no larger:
move first along its row to the second column, then down that column,
and use row and column standardness. Thus every box northwest of a
box with label less than $m$ also has label less than $m$. This proves
the Young diagram property. The original enumerating bijection
restricts to a bijection from the first $m$ labels onto these boxes,
and row and column inequalities are inherited. The inclusion of
initial label sets gives monotonicity. Summing box contents through
the restricted bijection gives the displayed equality; the next
prefix adds exactly the box carrying $m$, giving the difference.
\end{proof}

\begin{lemma}[Frozen-tail restriction of a polytabloid]
\label{lem:polytabloid-prefix-restriction}
\uses{lem:standard-prefix-diagram,lem:integral-specht-all-tableaux}
\lean{BostConjecture.Specht.prefixSlice,BostConjecture.Specht.prefixSlice_polytabloidInt,BostConjecture.Specht.prefixSlice_wordRepInt}\leanok
Let $u$ be the restriction of a standard tableau $t$ to its first $m$
labels. Extend a row word $w$ for $u$ to a row word $\widetilde w$
on all labels by keeping the rows of $t$ on the remaining labels.
Then $e_t(\widetilde w)=e_u(w)$. This coefficient restriction is
an integral linear map and intertwines permutations of the first
$m$ labels, extended by the identity on the other labels.
\end{lemma}
\begin{proof}
A column permutation contributing to $e_t(\widetilde w)$ must
preserve the row of each label at least $m$. Since it also preserves
that label's column, it fixes the label itself. The remaining action
is a column permutation of $u$. Conversely each column permutation
of $u$ extends uniquely by fixing the other labels, preserves its
sign, and contributes precisely its original term. This bijection
of the nonzero summands proves the formula. Restricting coordinates
is integral linear. Precomposing the word with a permutation of the
first $m$ labels commutes with extending it by the fixed tail, which
proves equivariance.
\end{proof}

\begin{lemma}[Support at the largest label]
\label{lem:polytabloid-last-support}
\uses{def:standard-tableau,lem:integral-specht-span}
\lean{BostConjecture.Specht.ColumnStandard.last_row_le_of_polytabloidInt_ne_zero}\leanok
If the largest label of a column-standard tableau $t$ lies in row
$r$, every word in the support of $e_t$ assigns that label a row
at most $r$.
\end{lemma}
\begin{proof}
A nonzero coefficient comes from a column permutation. The image of
the largest label is a label in its column. Column standardness
implies that its row is at most the row of the largest label.
\end{proof}

\begin{lemma}[Integral lifts and central action under restriction]
\label{lem:prefix-slice-central}
\uses{lem:polytabloid-prefix-restriction,lem:central-content-action}
\lean{BostConjecture.Specht.orbitLattice_le_map_prefixSlice,BostConjecture.Specht.prefixSlice_central_polytabloidInt}\leanok
Let $R_{t,m}$ be frozen-tail restriction and $u=t|_{\{0,\ldots,m-1\}}$.
Every vector of the integral lattice $L_u$ is $R_{t,m}(x)$ for some
$x\in L_t$. If $Z_m$ is the sum of transpositions of the first $m$
labels, acting on the original word space, then
\[
 R_{t,m}(Z_m e_t)=C_{\lambda_t^{(m)}}e_u.
\]
This identity concerns the restricted coefficients only.
\end{lemma}
\begin{proof}
Each generating orbit vector $g e_u$ lifts to $\widetilde g e_t$,
where $\widetilde g$ fixes the tail. These are integral vectors of
$L_t$, so their integral combinations give every required lift.
Equivariance and the polytabloid restriction formula identify
$R_{t,m}(Z_m e_t)$ with the smaller central operator applied to $e_u$.
The central content theorem on that smaller Specht module gives the
stated scalar.
\end{proof}

\begin{lemma}[Largest-label support filtration]
\label{lem:specht-last-row-filtration}
\uses{lem:polytabloid-last-support,lem:standard-polytabloid-independent,lem:real-standard-basis,lem:real-specht-action}
\lean{BostConjecture.Specht.lastRowSubspace,BostConjecture.Specht.standardRealBasis_repr_eq_zero_of_last_support,BostConjecture.Specht.mem_lastRowSubspace_iff,BostConjecture.Specht.realSpechtRep_mem_lastRowSubspace}\leanok
In the real Specht space of a diagram with $n+1$ boxes, let $F_r$
be the span of standard polytabloids whose largest label occupies
row at most $r$. Then
\[
 F_r=\{x:x(w)=0\text{ whenever }w(n)>r\}.
\]
Every permutation fixing the largest label preserves $F_r$.
\end{lemma}
\begin{proof}
The support lemma gives one inclusion. For the converse write $x$
in the standard basis. If a coefficient with largest-label row
exceeding $r$ is nonzero, choose among these the tableau $t$ with
lexicographically least row word. Evaluate the expansion at $w_t$.
A tableau with largest-label row at most $r$ contributes zero by
the support lemma. For every remaining nonzero coefficient, a
nonzero evaluation forces its leading row word to be no larger
than $w_t$, by column standardness. Minimality forces equality,
and row standardness then forces the tableau to be $t$.
The evaluation is therefore exactly the nonzero coefficient of $t$,
contradicting the assumed support. All coefficients outside $F_r$
vanish, proving the equality. A permutation fixing $n$ preserves
the final letter of each word and preserves the Specht space;
thus it preserves the displayed support condition and $F_r$.
\end{proof}

\begin{lemma}[Successive branching quotients]
\label{lem:specht-branching-quotient}
\uses{lem:specht-last-row-filtration,lem:polytabloid-prefix-restriction,lem:real-standard-basis}
\lean{BostConjecture.Specht.branchingMap_surjective,BostConjecture.Specht.ker_branchingMap,BostConjecture.Specht.branchingQuotientEquiv,BostConjecture.Specht.branchingQuotientEquiv_mk_apply,BostConjecture.Specht.branchingQuotientEquiv_mk_action}\leanok
Let $t$ be a standard tableau of $\mu$ on $n+1$ labels, let $r$ be
the row of its largest label, and let $\nu=\lambda_t^{(n)}$ be its
prefix diagram. The removed box is a removable corner. Let $F_{<r}$ be the span of
standard polytabloids with largest-label row strictly less than $r$.
Restriction to words assigning the largest label to row $r$ induces
an $\mathfrak S_n$-equivariant isomorphism
\[
 F_r/F_{<r}\simeq S^\nu_{\mathbb R}.
\]
On representatives it restricts word coefficients by assigning the
last label row $r$. In particular it sends the class of $e_t$ to
the polytabloid obtained by deleting its largest label.
\end{lemma}
\begin{proof}
Let $u$ be the prefix tableau of $t$. The largest label lies at the right end of
its row and bottom of its column, so its box is the specified corner.
Two row-standard tableaux with largest-label row $r$ have the same
last box. Consequently they differ by a permutation fixing the last
label. Frozen-tail restriction intertwines this permutation with its
restriction to the prefix. Thus every standard generator in row $r$
maps into $S^\nu_{\mathbb R}$, while every generator in a smaller row
maps to zero by the support bound. This proves the image inclusion.
Conversely, every permutation translate of $e_u$ is the restriction
of the same prefix permutation applied to $e_t$. These lifts remain
in $F_r$ by invariance. The translates span the smaller Specht space,
so restriction is onto.

For the kernel, suppose $x\in F_r$ has zero restriction. For each
standard tableau $v$ in row $r$, its leading row word is an extension
of a prefix word: use the prefix permutation carrying $t$ to $v$.
Thus $x(w_v)=0$. For standard tableaux in larger rows this evaluation
also vanishes by the support condition defining $F_r$. If any
standard-basis coefficient in row at least $r$ were nonzero, choose
one with lexicographically least leading word. Evaluating there,
the support bound eliminates the coefficients in earlier rows, and
leading-word triangularity eliminates all other nonzero coefficients
in rows at least $r$. The result is the chosen coefficient, a
contradiction. Hence only rows strictly less than $r$ occur, proving
that the kernel is $F_{<r}$. The first isomorphism theorem now gives
the asserted quotient isomorphism. Frozen-tail restriction gives
its formula on generators and its equivariance.
\end{proof}

\begin{lemma}[Last Murphy operator on a branching quotient]
\label{lem:last-murphy-branching}
\uses{lem:specht-branching-quotient,lem:standard-prefix-diagram,lem:central-content-action}
\lean{BostConjecture.Specht.prefixCentralOperator_sub_content_mem,BostConjecture.Specht.murphyOperator_last_sub_content_mem}\leanok
Let $t$ be a standard tableau on labels $0,\ldots,n$, set
$r=\operatorname{row}_t(n)$, and let $\nu=\lambda_t^{(n)}$.
Let $Z_n$ be the sum of transpositions of the first $n$ labels,
acting on the original real Specht space. For every $x\in F_r$,
\[
 Z_nx-C_\nu x\in F_{<r},\qquad
 J_nx-\operatorname{content}_t(n)x\in F_{<r}.
\]
\end{lemma}
\begin{proof}
Every summand of $Z_n$ fixes the last label, so it preserves $F_r$.
The branching restriction intertwines $Z_n$ with the full central
sum on the smaller Specht space. That sum is multiplication by
$C_\nu$. Thus $Z_nx-C_\nu x$ is in the kernel of restriction,
which is exactly $F_{<r}$ by the quotient theorem.
The full central sum is $Z_n+J_n$, since its remaining transpositions
are precisely those involving the last label. It acts by $C_\mu$.
The difference $C_\mu-C_\nu$ is the content of the removed box.
Consequently the second displayed difference is the negative of the
first, and belongs to the same subspace.
\end{proof}

\begin{lemma}[Extension at a fixed removable corner]
\label{lem:standard-corner-extension}
\uses{lem:standard-prefix-diagram,lem:polytabloid-prefix-restriction,lem:real-standard-basis}
\lean{BostConjecture.Specht.extendTableau_isStandard,BostConjecture.Specht.prefixSliceReal_extendTableau,BostConjecture.Specht.tableauColexLT_extendTableau_iff,BostConjecture.Specht.cornerLift_mem_lastRowSubspace,BostConjecture.Specht.prefixSliceReal_cornerLift,BostConjecture.Specht.cornerLift_mem_colexSubspace}\leanok
Let $t$ be standard on $n+1$ labels and let $\nu$ be its diagram
after deleting the largest label. Every standard tableau $u$ of
$\nu$ extends to a standard tableau $\widehat u$ of the original
diagram by putting label $n$ in the deleted box. Restriction of
$e_{\widehat u}$ at that last row is $e_u$. Colexicographic
comparison of two such extensions agrees with comparison of their
prefix row words.
The linear extension of this assignment is a right inverse to
restriction, with image in the chosen largest-label filtration
step, and it sends every earlier-prefix span into the corresponding
earlier-tableau span.
\end{lemma}
\begin{proof}
The removed box is at the right end of its row and the bottom of
its column. All row and column comparisons inside $\nu$ are
preserved. Every comparison involving the added box places its
largest label after the other label, so the extension is standard.
Choose the permutation of the first $n$ labels that carries the
prefix of $t$ to $u$. Its extension fixing $n$ carries $t$ to
$\widehat u$. Equivariance of frozen-tail restriction therefore
identifies the restricted polytabloid with $e_u$. The last row
agrees for both extensions, and their earlier row entries are
exactly those of the two prefixes. Their last differing label and
the row comparison at that label are consequently unchanged.
Extend the assignment linearly using the standard basis. The
restriction identity on basis vectors proves the right-inverse
identity on the whole smaller Specht space. Each image basis
vector has the prescribed final row; the order comparison just
proved gives the assertion on earlier spans.
\end{proof}

\begin{lemma}[Colexicographic triangular Murphy action]
\label{lem:prefix-content-triangular}
\uses{lem:specht-branching-quotient,lem:last-murphy-branching,lem:standard-corner-extension}
\lean{BostConjecture.Specht.murphyOperator_self_sub_content_mem_colex,BostConjecture.Specht.murphyOperator_sub_content_mem_colex}\leanok
Order standard tableaux by comparing their row words at the largest
label where they differ; the tableau with the smaller row there is
earlier. For every label $j$,
\[
 J_je_t-\operatorname{content}_t(j)e_t
 \in\operatorname{span}_{\mathbb R}\{e_s:s<t\}.
\]
This total order is sufficient for monic orthogonalization; the assertion
does not require triangularity for the stronger partial order comparing
every prefix in dominance order.
\end{lemma}
\begin{proof}
Induct on the number of labels. With no labels the assertion is empty.
For a diagram with $n+1$ boxes, the order first groups tableaux by
the row of label $n$. Inside a fixed group it is exactly the same
order on the tableaux obtained by deleting that label. For $j<n$,
the branching-quotient isomorphism and the induction hypothesis
show that $J_je_t-\operatorname{content}_t(j)e_t$ is a combination
of earlier tableaux in the same group, modulo $F_{<r}$. Every
tableau in $F_{<r}$ is earlier as well, proving the assertion.
For $j=n$, the last-Murphy branching lemma places
$J_ne_t-\operatorname{content}_t(n)e_t$ directly in $F_{<r}$.
Every tableau in that subspace is earlier, completing the induction.
\end{proof}

\begin{lemma}[Content eigenvalues of the monic orthogonal basis]
\label{lem:orthogonal-content-eigenvalues}
\uses{lem:orthogonal-standard-basis,lem:murphy-operators,lem:prefix-content-triangular}
\lean{BostConjecture.Specht.murphyOperator_mem_colexSubspace,BostConjecture.Specht.murphyOperator_orthogonalStandardBasis}\leanok
For the colexicographic monic orthogonal basis, every Murphy operator
acts diagonally:
\[
 J_j f_t=\operatorname{content}_t(j) f_t.
\]
\end{lemma}
\begin{proof}
Let $V_{<t}$ be the span of the earlier standard polytabloids.
Triangularity implies that $J_j$ preserves this span: the image of
each earlier generator is its content multiple plus still earlier
generators. Orthogonalization gives $f_t-e_t\in V_{<t}$ and
$f_t\perp V_{<t}$. Therefore
$z=J_j f_t-\operatorname{content}_t(j)f_t$ belongs to $V_{<t}$.
For every $y\in V_{<t}$, symmetry gives
$\langle z,y\rangle=\langle f_t,J_jy\rangle-
\operatorname{content}_t(j)\langle f_t,y\rangle=0$.
Taking $y=z$ proves $\|z\|^2=0$, hence $z=0$.
\end{proof}

\begin{lemma}[Content sequences distinguish standard tableaux]
\label{lem:standard-content-separation}
\uses{def:standard-tableau}
\lean{BostConjecture.Specht.IsStandard.eq_of_tableauContent_eq}\leanok
Two standard tableaux of the same shape with the same content at every
label are equal.
\end{lemma}
\begin{proof}
Induct on the label, assuming the boxes of all smaller labels agree.
The two boxes of the current label lie on the same diagonal, so they
are comparable in the northwest order. If distinct, the northwest box
has a strictly smaller label in the tableau where the current label
occupies the southeast box: labels are monotone along rows and columns.
The induction hypothesis then puts that smaller label into the same box
in the other tableau, contradicting its current label. Thus the two
boxes coincide, completing the induction.
\end{proof}

\begin{lemma}[Adjacent spectral support]
\label{lem:adjacent-spectral-support}
\uses{lem:standard-content-separation,lem:orthogonal-content-eigenvalues,lem:murphy-recurrence}
\lean{BostConjecture.Specht.eq_smul_orthogonalStandardBasis_of_eigenvalues,BostConjecture.Specht.murphyOperator_swap_residual,BostConjecture.Specht.exists_swap_orthogonalStandardBasis_eq}\leanok
Suppose $u=(i,i+1)t$ is standard and
$\Delta=\operatorname{content}_t(i+1)-\operatorname{content}_t(i)\ne0$.
Then $s f_t-\Delta^{-1}f_t$ is a scalar multiple of $f_u$, where
$s$ is the position transposition $(i,i+1)$.
\end{lemma}
\begin{proof}
For labels other than $i,i+1$, the Murphy operator commutes with $s$.
For larger labels this follows by permuting the summands of the Murphy
sum; for smaller labels every transposition in that sum has disjoint
support from $s$. The adjacent recurrence gives
$J_{i+1}s=sJ_i+1$ and $J_i s=sJ_{i+1}-1$.
Applying these identities to $f_t$ shows that the residual has the
content eigenvalues of $u$: the two adjacent eigenvalues are exchanged,
and all other eigenvalues are unchanged.
Expand the residual in the orthogonal basis. Symmetry of each Murphy
operator forces a nonzero coefficient to have all the same eigenvalues
as the residual. Content separation forces its index to be $u$.
\end{proof}

\begin{theorem}[Normalized adjacent action on the monic basis]
\label{lem:monic-seminormal}\uses{lem:orthogonal-standard-basis,lem:murphy-eigenvector-residual,lem:adjacent-spectral-support}
\lean{BostConjecture.Specht.realSpechtRep_swap_orthogonalStandardBasis,BostConjecture.Specht.norm_orthogonalStandardBasis_swap_sq}
\source{bost-specht-manuscript:page-0005,ryom-hansen-seminormal-denominators-v1:classical_formula}\leanok
Order standard tableaux colexicographically by their row words.
For a legal adjacent swap $u=(j,j+1)t$ with $t<u$, the monic
orthogonal basis satisfies
\[
 (j,j+1)f_t=\rho f_t+f_u,\qquad
 \|f_u\|^2=(1-\rho^2)\|f_t\|^2,\qquad
 \rho=\frac{1}{\operatorname{content}_t(j+1)-\operatorname{content}_t(j)}.
\]
These are the rational polytabloid normalization and left action of
\cite{ryom-hansen-seminormal}, Section 2, equation (1), after replacing
its index $i$ by $j+1$ and extending the rational inner product to
the real word realization.
\end{theorem}
\begin{proof}
The two contents differ: if they were equal, the legal swap would
preserve the entire content sequence and hence the tableau, contrary
to injectivity of its labeling. Spectral support gives
$sf_t=\rho f_t+c f_u$ and $sf_u\in\operatorname{span}(f_t,f_u)$.
Let $z=f_t-e_t$. It belongs to the earlier-tableau span $V_{<t}$,
which is perpendicular to both $f_t$ and $f_u$ since $t<u$.
Symmetry of $s$ gives $\langle f_u,sz\rangle=\langle sf_u,z\rangle=0$.
Literal relabeling of polytabloids gives $se_t=e_u$. Therefore
\[
 \langle f_u,sf_t\rangle=\langle f_u,e_u\rangle=\|f_u\|^2.
\]
The spectral-support formula makes the left side $c\|f_u\|^2$;
nonvanishing of $f_u$ gives $c=1$. Substituting the normalized action
in the adjacent residual norm identity proves the squared-norm ratio.
\end{proof}

\begin{lemma}[Initial column-reading vector]
\label{lem:column-reading-initial}
\uses{def:standard-tableau,lem:orthogonal-standard-basis,lem:rectangular-initial-norm}
\lean{BostConjecture.Specht.rectangleColumnTableau,BostConjecture.Specht.rectangleColumnTableau_isStandard,BostConjecture.Specht.rectangleColumnTableau_colexLT,BostConjecture.Specht.orthogonalStandardBasis_rectangleColumnTableau,BostConjecture.Specht.norm_orthogonalStandardBasis_rectangleColumnTableau_sq}\leanok
The column-reading tableau $t_0$ of the $a$ by $k$ rectangle is
standard and is first in colexicographic order. Consequently
$f_{t_0}=e_{t_0}$ and $\|f_{t_0}\|^2=(a!)^k$.
\end{lemma}
\begin{proof}
Compare any standard tableau $t$ with $t_0$, starting at the largest
label. Suppose all larger labels occupy their column-reading boxes,
and let $m$ be the first label placed differently. Its box in $t$
cannot be a box with larger column-reading label, because those are
already occupied. Thus its column is no greater than that of $m$
in $t_0$. If its row were also no greater, its box would lie strictly
northwest of the column-reading box of $m$. Standardness would force
the latter box to have a label larger than $m$, contradicting the
fixed larger-label boxes. Hence its row is greater. In particular
the largest differing row makes $t_0<t$. The column-reading tableau
itself is standard since its label in row $r$, column $q$ is $aq+r$.
Orthogonalization leaves the first vector unchanged, and its initial
squared norm is the column-group order $(a!)^k$.
\end{proof}

\begin{lemma}[Adjacent eigenvector residual]
\label{lem:murphy-eigenvector-residual}\uses{lem:murphy-recurrence,lem:real-specht-action}
\lean{BostConjecture.Specht.inner_swap_of_murphy_eigenvector,BostConjecture.Specht.inner_swap_residual_of_murphy_eigenvector,BostConjecture.Specht.norm_swap_residual_of_murphy_eigenvector}\leanok
Suppose $x\in V_t$, $J_i x=\alpha x$ and
$J_{i+1}x=\beta x$, with $\Delta=\beta-\alpha\ne0$.
For $s=P_{(i,i+1)}$ and $y=sx-\Delta^{-1}x$ one has
\[
 \langle x,sx\rangle=\|x\|^2/\Delta,\qquad
 \langle x,y\rangle=0,\qquad
 \|y\|^2=(1-\Delta^{-2})\|x\|^2.
\]
\end{lemma}
\begin{proof}
The recurrence gives $J_{i+1}sx=sJ_ix+x=\alpha sx+x$.
Taking its inner product with $x$ and using symmetry of $J_{i+1}$
gives $(\beta-\alpha)\langle x,sx\rangle=\|x\|^2$.
Division by $\Delta$ proves the first identity, and subtracting
$\Delta^{-1}\langle x,x\rangle$ proves the second. Since $s$
is an isometry, expansion of the squared norm gives
$\|y\|^2=\|x\|^2-2\Delta^{-2}\|x\|^2+
\Delta^{-2}\|x\|^2$, proving the last identity.
These statements assume the two eigenvalue equations; they do not
identify the eigenvalues of the orthogonalized standard polytabloids.
\end{proof}

\begin{lemma}[Legal adjacent paths from column reading]
\label{lem:standard-adjacent-path}
\uses{lem:column-reading-initial,def:standard-tableau}
\lean{BostConjecture.Specht.IsStandard.relabel_adjacent,BostConjecture.Specht.exists_rectangle_adjacent_descent,BostConjecture.Specht.rectangle_standard_adjacent_path}\leanok
Every standard tableau of a rectangle can be reached from column reading
by a finite sequence of legal adjacent swaps, each increasing colex order.
\end{lemma}
\begin{proof}
Let $t$ differ from column reading $t_0$ and form the permutation
$p=t_0^{-1}t$ of the labels. A strictly increasing permutation of a finite
chain is the identity, so some consecutive labels $i,i+1$ satisfy
$p(i)>p(i+1)$. Their boxes in $t$ must lie in different columns, with
$i$ in the later column: column reading orders columns first, and
standardness orders labels within each column. The box of $i$ is in a
strictly earlier row, since otherwise the box of $i+1$ is northwest of
it, contradicting standardness. Thus the two boxes have different rows
and columns. Swapping their consecutive labels preserves every row and
column comparison: only the relative order of these two labels changes,
and they share neither a row nor a column. The resulting tableau is
strictly earlier in colex order, since the largest changed label is
$i+1$ and its new row is earlier. Induction in the finite colex order
reaches $t_0$. Reverse the swaps to obtain the stated path.
\end{proof}

\begin{definition}[Prefix insertion product]
\label{def:prefix-norm-product}\uses{def:standard-tableau}
\lean{BostConjecture.Specht.prefixRowCount,BostConjecture.Specht.prefixAxialDistance,BostConjecture.Specht.prefixInsertionFactor,BostConjecture.Specht.prefixNormProduct}\leanok
For a standard enumerating tableau $t$, let $N_t(m,r)$ be the number of labels
strictly less than $m$ in row $r$. If label $m$ occupies $(s,q)$, put
\[
 D_t(m,r)=N_t(m,r)-q+s-r,\qquad
 P(t)=\prod_m\prod_{r<s}\frac{D_t(m,r)}{D_t(m,r)-1}.
\]
All row and column indices here start at zero.
\end{definition}

\begin{lemma}[Positive insertion denominators]
\label{lem:prefix-norm-positive}\uses{def:prefix-norm-product}
\lean{BostConjecture.Specht.prefixRowCount_self,BostConjecture.Specht.col_add_one_le_prefixRowCount,BostConjecture.Specht.two_le_prefixAxialDistance,BostConjecture.Specht.prefixNormProduct_pos}\leanok
For a standard tableau with label $m$ in $(s,q)$, one has
$N_t(m,s)=q$ and $N_t(m,r)\geq q+1$ for every $r<s$.
Consequently $D_t(m,r)\geq2$ for $r<s$ and $P(t)>0$.
\end{lemma}
\begin{proof}
The earlier labels of row $s$ occupy exactly its boxes in columns
$0,\ldots,q-1$, by row standardness. For $r<s$, the boxes
$(r,0),\ldots,(r,q)$ all belong to the diagram because $(s,q)$ does.
Each is strictly northwest of $(s,q)$, so its label is strictly
smaller than $m$. These $q+1$ distinct boxes give the lower bound.
Since $s-r\geq1$, the displayed formula for $D_t(m,r)$ is at least
two. Every numerator and denominator in the finite product is positive.
\end{proof}

\begin{lemma}[Adjacent ratio of the prefix product]
\label{lem:prefix-norm-swap}\uses{lem:prefix-norm-positive}
\lean{BostConjecture.Specht.prefixNormProduct_relabel_adjacent,BostConjecture.Specht.prefixAxialDistance_eq_content_sub,BostConjecture.Specht.two_le_prefixAxialDistance}\leanok
Let $t$ be standard, with consecutive labels $m,m+1$ in rows $s,r$
where $r<s$, and let $u=(m,m+1)t$. Put
$\Delta=\operatorname{content}_t(m+1)-\operatorname{content}_t(m)$.
Then $\Delta\geq2$ and
\[
 P(u)=(1-\Delta^{-2})P(t).
\]
\end{lemma}
\begin{proof}
Only the prefix between the two labels changes. The upper insertion
has identical factors in both tableaux because the lower row contributes
to none of its upper-row counts. For the lower insertion, only the factor
belonging to the upper row $r$ changes. Its distance increases from
$D$ to $D+1$. The ratio of these factors is
\[
 \frac{(D+1)/D}{D/(D-1)}=1-D^{-2}.
\]
The upper row's count before the lower insertion is the column of label
$m+1$: adding the lower box does not change that count. Hence $D$ equals
the displayed content difference $\Delta$, and positivity gives $D\geq2$.
All other factors are unchanged, proving the product identity.
\end{proof}

\begin{lemma}[Column-reading prefix product]
\label{lem:prefix-norm-initial}\uses{def:prefix-norm-product,lem:column-reading-initial}
\lean{BostConjecture.Specht.prefixNormProduct_rectangleColumnTableau}\leanok
For column reading $t_0$ of the $a$ by $k$ rectangle,
$P(t_0)=(a!)^k$.
\end{lemma}
\begin{proof}
For a box in row $s$, column $q$, each row $r<s$ has exactly $q+1$
earlier entries, namely its boxes in columns $0,\ldots,q$.
Thus $D_{t_0}(m,r)=s-r+1$ and the insertion product telescopes:
$\prod_{r<s}(s-r+1)/(s-r)=s+1$.
Multiplication down each column gives $a!$, and there are $k$ columns.
\end{proof}

\begin{lemma}[Seminormal norms]
\label{lem:seminormal-norms}\uses{lem:monic-seminormal,lem:column-reading-initial,lem:standard-adjacent-path,lem:prefix-norm-swap,lem:prefix-norm-initial,lem:orthogonal-standard-basis}
\lean{BostConjecture.Specht.prefixRowCount_eq_rowLen_prefixShape,BostConjecture.Specht.norm_orthogonalStandardBasis_sq_eq_prefixNormProduct,BostConjecture.Specht.det_standard_gram_eq_prod_prefixNormProduct}
\source{bost-specht-manuscript:page-0005,bost-specht-manuscript:page-0006}\leanok
For a standard tableau $t$, let $\lambda^{(j)}$ be its prefix shape
after $j$ entries, and let $r_j$ be the row of entry $j$. Set
$\delta_{i,r_j}(j)=\lambda_i^{(j-1)}-\lambda_{r_j}^{(j-1)}+r_j-i$.
Then each relevant $\delta$ is at least two, and
\[
 \gamma_t:=\|f_t\|^2
 =\prod_{j=1}^d\prod_{i<r_j}\frac{\delta_{i,r_j}(j)}{\delta_{i,r_j}(j)-1},
 \qquad
 \det\Gram((e_t)_{t\in\SYT(k^a)})=\prod_t\gamma_t.
\]
\end{lemma}
\begin{proof}
Earlier labels in a fixed row are exactly the boxes in that row of the
prefix diagram. Thus $N_t(j,i)=\lambda_i^{(j-1)}$, and the current
row's earlier-label count is the column of the inserted box. The
explicit product $P(t)$ is therefore the product displayed here.
Every distance is at least two by positivity of insertion denominators.

The actual squared norm and $P(t)$ both equal $(a!)^k$ at column
reading. Every standard tableau is reached by a finite increasing
legal adjacent path. At each step, the normalized adjacent action
and the prefix-product ratio multiply the two quantities by exactly
the same factor $1-\Delta^{-2}$. Induction along the path proves
$\|f_t\|^2=P(t)$ for every standard tableau. Finally, the monic
orthogonal change of basis preserves the Gram determinant, and the
orthogonal Gram matrix has diagonal entries $\|f_t\|^2$.
\end{proof}

\begin{lemma}[Seminormal upper bound for orbit covolume]
\label{lem:seminormal-covolume}
\uses{lem:seminormal-norms,lem:real-standard-basis,lem:standard-gram-covolume,lem:rectangular-covolume-nonneg}
\lean{BostConjecture.Specht.two_mul_log_covolume_le_sum_log_prefixNormProduct,BostConjecture.Specht.two_mul_rectangleLogCovolume_le_average_log_prefixNormProduct}
\source{bost-specht-manuscript:page-0005,bost-specht-manuscript:page-0006}\leanok
For the all-polytabloid lattice,
\[
 0\leq2c_a(d)\leq\mathbb E_{t\in\SYT(k^a)}\log\gamma_t.
\]
\end{lemma}
\begin{proof}
The standard polytabloids are integral vectors forming a real basis,
so the covolume bound gives $0\leq2\log\covol L_{a,d}\leq
\log\det\Gram((e_t))$. Each $\gamma_t$ is positive. Taking the
logarithm of the seminormal Gram identity gives
$\log\det\Gram((e_t))=\sum_t\log\gamma_t$. Divide by the
positive number of standard tableaux, which equals $f_{a,d}$.
\end{proof}

\begin{lemma}[Tableau enumeration in shifted coordinates]
\label{lem:tableau-count}
\source{bost-specht-manuscript:page-0004,bost-specht-manuscript:page-0006}\notready
For a partition $\nu$ of $t$ with at most $a$ rows, set
$b_i=\nu_i+a-i$, $1\leq i\leq a$, and
$V(b)=\prod_{i<j}(b_i-b_j)$. The number of standard tableaux is
\[
 f^\nu=\frac{t!V(b)}{\prod_i b_i!}.
\]
Equivalently it is $t!$ divided by the product of all hook lengths.
\end{lemma}
\begin{proof}
We derive the shifted-coordinate form by corner recursion, giving a
direct proof of the enumeration formula used in the manuscript.
Put $c=a(a-1)/2$, so $t=\sum_i b_i-c$. For distinct $b_i$, expand
\[
 R(z)=\prod_i\frac{z-b_i-1}{z-b_i}
 =1-\sum_i\frac{p_i}{z-b_i},\qquad
 p_i=\prod_{j\ne i}\frac{b_i-b_j-1}{b_i-b_j}.
\]
The residues prove this identity: the difference has no poles and
tends to zero at infinity. The coefficient of $z^{-2}$ in the
product is $c-\sum_i b_i=-t$; in the partial fraction expression
it is $-\sum_i b_ip_i$. Hence $\sum_i b_ip_i=t$.
For $t>0$ and a legal decrease of $b_i$ by one, the ratio of the
proposed new count to the proposed old count is $b_ip_i/t$.
An illegal decrease contributes zero: either $b_i=0$, or
$b_i-b_{i+1}=1$, which makes the corresponding factor of $p_i$ zero.
Legal decreases are precisely removal of a corner box. The identity
therefore says the proposed count is the sum of its counts after
all possible corner removals. Standard tableaux satisfy the same
recurrence, since their largest entry lies in a unique removable
corner. The empty partition has $b_i=a-i$ and proposed count one,
since $V(a-1,\ldots,0)=\prod_{j=1}^{a-1}j!$.
Induction on $t$ proves the count.

For the hook version, the hook lengths in row $i$ are exactly the
integers $1,\ldots,b_i$ with $b_i-b_j$ omitted for $j>i$.
Indeed, for columns satisfying $\nu_{r+1}<q\leq\nu_r$ with $r\geq i$,
the hook length is $\nu_i-q+r-i+1$; the consecutive intervals of
these values have exactly those omitted endpoints. Multiplying row
products gives the hook product $\prod_i b_i!/V(b)$, proving the
equivalent formula.
\end{proof}

\begin{lemma}[Exact prefix distribution]
\label{lem:prefix-distribution}\uses{lem:tableau-count}
\source{bost-specht-manuscript:page-0006}\notready
Fix $a\geq2$, $k\geq1$, $d=ak$, $c=a(a-1)/2$, $M=k+a-1$,
and $V_0=\prod_{j=1}^{a-1}j!$. For a uniform standard tableau of
$(k^a)$, write $\nu$ for the prefix after $t$ entries and put
$b_i=\nu_i+a-i$, $A=t+c$. Its distribution is obtained by weighting
the labelled multivariate hypergeometric distribution
\[
 B(b)=\frac{\prod_i\binom M{b_i}}{\binom{aM}A},\qquad
 0\leq b_i\leq M,\quad \sum_i b_i=A,
\]
by $V(b)^2$ and sorting into decreasing order. The normalizing factor is
\[
 Z_{k,t}:=\mathbb E_B V(b)^2
 =a!V_0\prod_{j=1}^{a-1}(k+j)^j
   \frac{\binom dt}{\binom{d+2c}{t+c}}.
\]
\end{lemma}
\begin{proof}
The exact prefix probability is $f^\nu f^{\nu^c}/f^{(k^a)}$, where
$\nu^c=(k-\nu_a,\ldots,k-\nu_1)$. The first $t$ entries fill $\nu$
arbitrarily as a standard tableau. Rotation by $180$ degrees and
reversal of labels identifies standard fillings of the remaining
skew diagram with those of $\nu^c$, proving that probability.
The shifted coordinates of the complement are $M-b_{a+1-i}$, with
the same positive Vandermonde product. The enumeration formula gives
\[
 f^\nu=\frac{t!V(b)}{\prod_i b_i!},\quad
 f^{\nu^c}=\frac{(d-t)!V(b)}{\prod_i(M-b_i)!},\quad
 f^{(k^a)}=\frac{d!V_0}{\prod_{j=0}^{a-1}(k+j)!}.
\]
Their ratio is a constant times $V(b)^2B(b)$. Every decreasing tuple
has $a!$ labelled permutations, and repeated coordinates contribute
zero Vandermonde weight. Summing the probabilities therefore gives
\[
 Z_{k,t}=a!V_0\frac{(M!)^a}{\prod_{j=0}^{a-1}(k+j)!}
             \frac{\binom dt}{\binom{aM}{t+c}}.
\]
Use $aM=d+2c$ and
$(M!)^a/\prod_{j=0}^{a-1}(k+j)!=\prod_{j=1}^{a-1}(k+j)^j$
to obtain the stated expression.
\end{proof}

\begin{lemma}[Uniform separation of bulk prefix coordinates]
\label{lem:prefix-separation}\uses{lem:prefix-distribution}
\source{bost-specht-manuscript:page-0007}\notready
For fixed $a\geq2$, $0<\varepsilon<1/2$, and $L>0$,
\[
 \sup_{\varepsilon d\leq t\leq(1-\varepsilon)d}
 \mathbb P\left(\min_{i<j}|b_i-b_j|\leq L\right)\longrightarrow0
 \qquad(k\to\infty).
\]
\end{lemma}
\begin{proof}
The exact normalization satisfies
\[
 \frac{\binom dt}{\binom{d+2c}{t+c}}
 =\frac{\prod_{j=1}^c(t+j)(d-t+j)}{\prod_{j=1}^{2c}(d+j)}.
\]
For the stated range of $t$, every numerator factor is at least
$\varepsilon d$, while every denominator factor is at most
$d+2c=a(k+a-1)\leq ad$. Thus
$Z_{k,t}\geq a!V_0(\varepsilon/a)^{2c}k^c$.

Under the unweighted hypergeometric law, the moment-generating
function of $b_i$ is the normalized $A$th elementary symmetric
polynomial of $aM$ positive numbers, of which $M$ equal $e^\theta$
and the others equal one. For positive numbers with fixed sum, this
normalized polynomial is at most the $A$th power of their mean.
To prove that inequality, replace a largest and a smallest entry by
their mean. With their sum fixed, the symmetric polynomial is affine
in their product with nonnegative coefficient, so does not decrease.
The sum of squared deviations from the mean drops by half their
squared difference, at least the previous sum of squared deviations
divided by $2aM$. Repeating converges to the constant tuple; continuity
proves the claim. Consequently
\[
 \mathbb E_B e^{\theta b_i}\leq(1-1/a+e^\theta/a)^A.
\]
For $p=1/a$, the second derivative of
$\log(1-p+pe^\theta)$ is at most $1/4$, so subtracting its tangent
at zero yields
$\mathbb E_Be^{\theta(b_i-A/a)}\leq e^{A\theta^2/8}$.
Markov's inequality with $\theta=4u/A$ and $-4u/A$ gives
$\mathbb P_B(|b_i-A/a|>u)\leq2e^{-2u^2/A}$.

Put $\eta=1/(4c)$ and $u=k^{1/2+\eta}$. Since $A\leq aM\leq a^2k$,
the union bound gives
$\mathbb P_B(\max_i|b_i-A/a|>u)\leq2a e^{-2k^{2\eta}/a^2}$.
On the complementary event all pairwise differences are at most
$2u$. If one difference is at most $L$, then
$V(b)^2\leq L^2(2u)^{2(c-1)}$.
Dividing by the lower bound for $Z$ bounds the weighted probability
of this event by a constant times
$k^{-1+2\eta(c-1)}=k^{-1/2-1/(2c)}$, tending to zero. On the tail
event use $V(b)^2\leq M^{2c}$; its weighted probability is at most a
constant times $k^c e^{-2k^{2\eta}/a^2}$, also tending to zero.
The constants are independent of $t$ in the stated range. Both events
are invariant under sorting the coordinates, so the same bound
holds for the prefix distribution.
\end{proof}

\begin{lemma}[The one-row case]
\label{lem:one-row-covolume}\uses{def:rectangular-specht,lem:rectangular-initial-norm,lem:integral-specht-covolume-nonneg}
\lean{BostConjecture.Specht.oneRowIntegralEquiv,BostConjecture.Specht.oneRowIntegralBasis,BostConjecture.Specht.oneRow_covolume_eq_one,BostConjecture.Specht.rectangleLogCovolume_one,BostConjecture.Specht.tendsto_rectangleLogCovolume_one_div}
\source{bost-specht-manuscript:page-0004}\leanok
The one-row lattice is isomorphic to $\mathbb Z$ with a unit-length
generator. Thus $\covol L_{1,k}=1$, $c_1(k)=0$ and
$c_1(k)/k\to0$.
\end{lemma}
\begin{proof}
There is exactly one row word. Evaluation at it sends an integral
lattice vector to its sole coordinate and has inverse
$z\mapsto z e_t$, since the row-word coefficient of $e_t$ is one.
Thus $e_t$ is an integral basis. Its squared norm is $(1!)^k=1$,
so the one-by-one Gram determinant gives covolume one. Taking its
logarithm gives zero, and the sequence of normalized logarithmic
covolumes is identically zero.
\end{proof}

\begin{theorem}[Sublinear rectangular Specht covolume]
\label{thm:specht-sublinear-covolume}
\uses{def:rectangular-specht,lem:rectangular-covolume-nonneg,lem:one-row-covolume,lem:seminormal-norms,lem:seminormal-covolume,lem:prefix-separation}
\source{bost-specht-manuscript:page-0008}\notready
For every fixed positive integer $a$,
$c_a(ak)\geq0$ and $c_a(ak)/(ak)\to0$ as $k\to\infty$.
\end{theorem}
\begin{proof}
The integral Gram determinant lemma gives $c_a(d)\geq0$, since the
lattice has positive rank.
For $a=1$, the lattice has one orthonormal generator, so $c_1(d)=0$.
For $a\geq2$, each insertion contributes a number between zero and
$(a-1)\log2$ to $\log\gamma_t$. If all shifted coordinate gaps of
the preceding prefix exceed $L\geq2$, each factor contributes at
most $1/(L-1)$, using $\log(1+x)\leq x$. Let $p_{k,\varepsilon,L}$
be the supremum of the bad-prefix probabilities in the preceding
lemma. There are at most $2\varepsilon d+2$ insertion times whose
preceding prefix lies outside the bulk. Linearity of expectation,
with the bound valid for every next legal row, gives
\[
 0\leq\frac{\mathbb E\log\gamma_t}{d}
 \leq(2\varepsilon+2/d)(a-1)\log2
    +\frac{a-1}{L-1}+(a-1)\log2\,p_{k,\varepsilon,L}.
\]
For fixed $\varepsilon,L$ the last term tends to zero. Taking the
upper limit in $k$, then taking $L\to\infty$ and
$\varepsilon\to0$, proves $\mathbb E\log\gamma_t/d\to0$.
The inequalities $0\leq2c_a(d)/d\leq\mathbb E\log\gamma_t/d$
then give $c_a(d)/d\to0$ by squeezing.
\end{proof}
